\section{\texorpdfstring{\textcolor{blue}{Point 7: the group, or the gliders around it?}}{Point 7: the group, or the gliders around it?}}\label{sec:intergroup}

Group time progresses straighter than solo flight at the long lags. A group is a strict object: its members stay in contact long enough, and gliders a few kilometres ahead are not members, yet the pilot can see them. How much of the group effect belongs to the group itself (intra-group), and how much to the presence of other gliders, in other groups or alone (inter-group)? Only gliders ahead can show the way, so we also ask whether gliders ahead matter more than gliders behind.

\paragraph{Positions and outsiders.} The positions are those of the archive index: one position per minute of every flight declared to the federation, paragliders and hang gliders, on the sphere. An \emph{outsider} of a pilot at a minute is any other flight of the index at that minute that is not in the pilot's group then. For a solo pilot every other flight is an outsider. The minute is the resolution: distances below a few hundred metres and delays below a minute are not resolved.

\paragraph{Direction, ahead and behind.} The direction of motion of a pilot at the minute $t$ is the unit vector
\[
  \mathbf{u}(t)=\frac{\Pi\bigl[\mathbf{x}(t+T/2)-\mathbf{x}(t-T/2)\bigr]}{\bigl\|\Pi\bigl[\mathbf{x}(t+T/2)-\mathbf{x}(t-T/2)\bigr]\bigr\|},\qquad T=\IgTDir~\text{min},
\]
where $\mathbf{x}$ is the position and $\Pi$ projects on the horizontal plane at the pilot. The window is long enough to average the circling. We leave $\mathbf{u}$ undefined when the projected displacement is shorter than \IgMinDisp~m. While the pilot is in a group, the displacement is the mean over the members of the group, that is the displacement of its centroid.
An outsider $j$ is \emph{ahead} of the pilot $i$ when $(\mathbf{x}_j-\mathbf{x}_i)\cdot\mathbf{u}_i>0$ and behind otherwise. The gap ahead $d_a(t)$ is the horizontal distance to the nearest outsider ahead, the gap behind $d_b(t)$ the same behind; both are infinite when no outsider lies within \IgRMax~km.

\paragraph{The geometry of a group.} At every minute of a group of at least two members with a direction, let $\mathbf{c}$ be the centroid of the members and $\mathbf{w}$ the horizontal unit vector perpendicular to $\mathbf{u}$. For the member $i$ we write $s_i=(\mathbf{x}_i-\mathbf{c})\cdot\mathbf{u}$ and $p_i=(\mathbf{x}_i-\mathbf{c})\cdot\mathbf{w}$. The length and the width of the group are
\[
  L_\parallel=\max_i s_i-\min_i s_i,\qquad L_\perp=\max_i p_i-\min_i p_i,
\]
and its leader is the member with the largest $s_i$. The gap of the group is $d_a$ measured from the leader, with the members left out, split by the kind of outsider (a member of another group, or a solo glider); we also record the cross-track offset of the nearest outsider ahead. The groups are those with at least one open flight: \IgGroups{} groups and \IgGroupMinutes{} group minutes.

\begin{figure}
\centering
\includegraphics[width=\textwidth]{open_min0.1/intergroup_geometry.pdf}
\caption{Left: the median (dot) and the range from the 10th to the 90th percentile of the length and the width of the groups and of the gap ahead of the leader, on one distance axis. The gaps are those of the minutes that have an outsider ahead within \IgRMax~km; the share of minutes with none is written beside each. Right: the share of the time with an outsider ahead within $R$: group minutes seen from the leader, the pieces of $\G$ seen from the pilot, and the pieces of $\A$ and of $\Sl$ seen from the pilot.}
\label{fig:ig-geometry}
\end{figure}

\paragraph{What the groups look like.} The median group is \IgLenMed~m long (10th to 90th percentile \IgLenLow{} to \IgLenHigh~m) and \IgWidMed~m wide (\IgWidLow{} to \IgWidHigh~m). When an outsider is ahead of the leader within \IgRMax~km, the median gap is \IgGapMed~km (\IgGapLow{} to \IgGapHigh~km); in \IgGapNone\% of the group minutes there is none. For another group the median gap is \IgGapOtherMed~km and \IgGapOtherNone\% of the minutes have none. An outsider is ahead of the leader within \IgRadii~km during \IgShareGroupOne, \IgShareGroupTwo, \IgShareGroupFive{} and \IgShareGroupTen\% of the group time, and a member of another group during \IgShareOtherOne, \IgShareOtherTwo, \IgShareOtherFive{} and \IgShareOtherTen\%. The pieces of $\Sl$ have an outsider ahead during \IgShareSOne, \IgShareSTwo, \IgShareSFive{} and \IgShareSTen\% of their time, the pieces of $\A$ during \IgShareAOne, \IgShareATwo, \IgShareAFive{} and \IgShareATen\%, so a group is more often near others than a solo pilot, and the solo time of a group flight is the most crowded of the three. We have not compared these shares with a chance level: the time-shift surrogates of the grouping do not carry the positions of the gliders around.

\paragraph{The neighbourhood of a piece.} The units are the pieces of the raw design (dataset 1000, at least \RawShortestOne~s), with $\Sl$ and $\A$ pooled as solo time and $\G$ as group time. For a piece with $n$ minutes that have a direction,
\[
  f_a(R)=\frac1n\sum_t\mathbf{1}\bigl[d_a(t)\le R\bigr],\qquad f_b(R)=\frac1n\sum_t\mathbf{1}\bigl[d_b(t)\le R\bigr].
\]
A piece enters when the direction is defined in at least \IgMinDefined\% of its minutes: \IgPiecesUsed{} of the \IgPieces{} pieces. With $R=\IgRMain$~km in the main line, and $F_L=\IgFLow$, $F_H=\IgFHigh$, the class of a piece is
\begin{itemize}
  \item ALONE: $f_a<F_L$ and $f_b<F_L$;
  \item AHEAD: $f_a\ge F_H$;
  \item BEHIND: $f_b\ge F_H$ and $f_a<F_L$, so the others are behind and the pilot leads them;
  \item MIXED: the rest, which we count and show but do not interpret.
\end{itemize}
The classes at \IgRadLow{} and \IgRadHigh~km are in Appendix~\ref{app:intergroup}.

\paragraph{The effects.} Let $H(\text{state},\text{class})$ stand for $\Hk$ of the pieces of a state (solo or group) and a class. We take six contrasts:
\begin{align*}
  \Delta_\text{intra} &= H(\text{group},\text{ALONE})-H(\text{solo},\text{ALONE}), & \Delta_\text{inter} &= H(\text{solo},\text{AHEAD})-H(\text{solo},\text{ALONE}),\\
  \Delta_\text{crowd} &= H(\text{solo},\text{BEHIND})-H(\text{solo},\text{ALONE}), & \Delta_\text{info,solo} &= H(\text{solo},\text{AHEAD})-H(\text{solo},\text{BEHIND}),\\
  \Delta_\text{info,group} &= H(\text{group},\text{AHEAD})-H(\text{group},\text{BEHIND}), & \Delta_\text{both} &= H(\text{group},\text{AHEAD})-H(\text{solo},\text{ALONE}).
\end{align*}
The first is a group with nobody around against a pilot with nobody around. The second is a pilot with gliders ahead and no group. The third is the control: gliders behind cannot show the way, so it measures what a frequented route gives without information. The fourth and the fifth separate information from the frequented route, solo and in a group.
We estimate them in two lines. In the raw line, $H$ of a cell is $\Hk$ of the ensemble MSD of its pieces, with equal weight per flight, as in Section~\ref{sec:raw}; the contrast is the difference of the two. In the within-flight line, in analogy with the comparison of the same flight of Section~\ref{sec:decomposition}, we keep the flights that have pieces in both cells $P$ and $Q$ and take the local slope, over the lags that give $\Hk$, of the mean over those flights of $\log_{10}\MSD_P-\log_{10}\MSD_Q$. For the phase statistics the raw line is the estimator that averages first and takes the ratio after (unit the flight), and the within-flight line the mean over the flights with both cells of the difference of the ratio of sums of the flight. The bands are the two-way bands of the comparisons, from the same draws for every number. We show every cell with its counts and do not merge classes to rescue a thin one.

\begin{table}
\centering\small
\begin{tabular}{@{}l l r r r l@{}}
\toprule
state & class & pieces & flights & pilots & $\Hk$ \\
\midrule
solo & alone & \IgCellFiveSoloAlonePieces & \IgCellFiveSoloAloneFlights & \IgCellFiveSoloAlonePilots & \IgCellFiveSoloAloneH{} \IgCellFiveSoloAloneHBand \\
solo & ahead & \IgCellFiveSoloAheadPieces & \IgCellFiveSoloAheadFlights & \IgCellFiveSoloAheadPilots & \IgCellFiveSoloAheadH{} \IgCellFiveSoloAheadHBand \\
solo & behind & \IgCellFiveSoloBehindPieces & \IgCellFiveSoloBehindFlights & \IgCellFiveSoloBehindPilots & \IgCellFiveSoloBehindH{} \IgCellFiveSoloBehindHBand \\
solo & mixed & \IgCellFiveSoloMixedPieces & \IgCellFiveSoloMixedFlights & \IgCellFiveSoloMixedPilots & \IgCellFiveSoloMixedH{} \IgCellFiveSoloMixedHBand \\
group & alone & \IgCellFiveGroupAlonePieces & \IgCellFiveGroupAloneFlights & \IgCellFiveGroupAlonePilots & \IgCellFiveGroupAloneH{} \IgCellFiveGroupAloneHBand \\
group & ahead & \IgCellFiveGroupAheadPieces & \IgCellFiveGroupAheadFlights & \IgCellFiveGroupAheadPilots & \IgCellFiveGroupAheadH{} \IgCellFiveGroupAheadHBand \\
group & behind & \IgCellFiveGroupBehindPieces & \IgCellFiveGroupBehindFlights & \IgCellFiveGroupBehindPilots & \IgCellFiveGroupBehindH{} \IgCellFiveGroupBehindHBand \\
group & mixed & \IgCellFiveGroupMixedPieces & \IgCellFiveGroupMixedFlights & \IgCellFiveGroupMixedPilots & \IgCellFiveGroupMixedH{} \IgCellFiveGroupMixedHBand \\
\midrule
\multicolumn{6}{@{}l}{$\Delta\Hk$: raw \quad within the flight (flights with both cells)} \\
\multicolumn{2}{@{}l}{intra-group: (group, alone) $-$ (solo, alone)} & & \multicolumn{2}{l}{\IgConFiveIntraRaw{} \IgConFiveIntraRawBand} & \IgConFiveIntraFlight{} \IgConFiveIntraFlightBand{} (\IgConFiveIntraFlights) \\
\multicolumn{2}{@{}l}{inter-group: (solo, ahead) $-$ (solo, alone)} & & \multicolumn{2}{l}{\IgConFiveInterRaw{} \IgConFiveInterRawBand} & \IgConFiveInterFlight{} \IgConFiveInterFlightBand{} (\IgConFiveInterFlights) \\
\multicolumn{2}{@{}l}{crowd control: (solo, behind) $-$ (solo, alone)} & & \multicolumn{2}{l}{\IgConFiveCrowdRaw{} \IgConFiveCrowdRawBand} & \IgConFiveCrowdFlight{} \IgConFiveCrowdFlightBand{} (\IgConFiveCrowdFlights) \\
\multicolumn{2}{@{}l}{information, solo: (solo, ahead) $-$ (solo, behind)} & & \multicolumn{2}{l}{\IgConFiveInfoSoloRaw{} \IgConFiveInfoSoloRawBand} & \IgConFiveInfoSoloFlight{} \IgConFiveInfoSoloFlightBand{} (\IgConFiveInfoSoloFlights) \\
\multicolumn{2}{@{}l}{information, group: (group, ahead) $-$ (group, behind)} & & \multicolumn{2}{l}{\IgConFiveInfoGroupRaw{} \IgConFiveInfoGroupRawBand} & \IgConFiveInfoGroupFlight{} \IgConFiveInfoGroupFlightBand{} (\IgConFiveInfoGroupFlights) \\
\multicolumn{2}{@{}l}{both: (group, ahead) $-$ (solo, alone)} & & \multicolumn{2}{l}{\IgConFiveBothRaw{} \IgConFiveBothRawBand} & \IgConFiveBothFlight{} \IgConFiveBothFlightBand{} (\IgConFiveBothFlights) \\
\bottomrule
\end{tabular}

\caption{The cells of the state and the class at $R=\IgRMain$~km, with the pieces, the flights and the pilots of each, $\Hk$ with its 90\% band, and the six contrasts of $\Hk$, raw and within the flight (flights with both cells in brackets). BEHIND is the thinnest cell and MIXED is not interpreted. At 2 and \IgRadHigh~km see Appendix~\ref{app:intergroup}.}
\label{tab:intergroup}
\end{table}

\begin{figure}
\centering
\includegraphics[width=\textwidth]{open_min0.1/intergroup_effects.pdf}
\caption{The six contrasts at $R=\IgRMain$~km for $\Hk$ and the four phase statistics: raw (filled, blue) and within the flight (open, orange), with 90\% bands. From top to bottom: intra-group, inter-group, crowd control, information in solo flight, information in a group, both.}
\label{fig:ig-effects}
\end{figure}

\paragraph{The intra-group effect.} A group with nobody around is a little straighter than a solo pilot with nobody around: $\Delta_\text{intra}$ is \IgConFiveIntraRaw{} \IgConFiveIntraRawBand{} in the raw line and \IgConFiveIntraFlight{} \IgConFiveIntraFlightBand{} within the flight, on \IgConFiveIntraFlights{} flights. The raw gap $\G-\Sl$ of Section~\ref{sec:raw} is \RawGapHK, and the within-flight band includes 0. The group also climbs more (\IgPhFiveIntraClimbRaw{} points of time in climb, raw) and glides flatter (\IgPhFiveIntraSinkRaw~m/km).

\paragraph{The inter-group effect.} A solo pilot with gliders ahead is less straight than one with nobody around in the raw line: $\Delta_\text{inter}$ is \IgConFiveInterRaw{} \IgConFiveInterRawBand. Within the flight it is \IgConFiveInterFlight{} \IgConFiveInterFlightBand{} on \IgConFiveInterFlights{} flights, and the band includes 0. The two lines do not agree, so we read the raw difference as a matter of who flies where and when. The phase statistics point the same way as the second line: within the flight a pilot with gliders ahead spends \IgPhFiveInterClimbFlight{} \IgPhFiveInterClimbFlightBand{} points more of the time in climb, finds a climb after \IgPhFiveInterFoundFlight{} \IgPhFiveInterFoundFlightBand{} points more of its searches and loses \IgPhFiveInterSinkFlight{} \IgPhFiveInterSinkFlightBand~m/km, while the raw time in climb is \IgPhFiveInterClimbRaw{} points.
The control does not behave like the inter-group effect for $\Hk$: $\Delta_\text{crowd}$ is \IgConFiveCrowdRaw{} \IgConFiveCrowdRawBand{} raw and \IgConFiveCrowdFlight{} \IgConFiveCrowdFlightBand{} within the flight, on \IgConFiveCrowdFlights{} flights. For the time in climb the control gives \IgPhFiveCrowdClimbFlight{} \IgPhFiveCrowdClimbFlightBand{} within the flight, close to the inter-group value.

\paragraph{Ahead against behind.} The contrast $\Delta_\text{info,solo}$ is \IgConFiveInfoSoloRaw{} \IgConFiveInfoSoloRawBand{} raw and \IgConFiveInfoSoloFlight{} \IgConFiveInfoSoloFlightBand{} within the flight (\IgConFiveInfoSoloFlights{} flights): a pilot with gliders ahead is less straight than one with gliders behind, and both bands exclude 0. Inside a group $\Delta_\text{info,group}$ is \IgConFiveInfoGroupRaw{} \IgConFiveInfoGroupRawBand{} raw and \IgConFiveInfoGroupFlight{} \IgConFiveInfoGroupFlightBand{} within the flight (\IgConFiveInfoGroupFlights{} flights). The sign is the opposite of the one that information from the gliders ahead would give. A pilot with gliders only behind leads them and is also the pilot who flies straighter in the raw line. Having both groups and gliders ahead, $\Delta_\text{both}$ is \IgConFiveBothRaw{} \IgConFiveBothRawBand{} raw and \IgConFiveBothFlight{} \IgConFiveBothFlightBand{} within the flight.

\begin{figure}
\centering
\includegraphics[width=\textwidth]{open_min0.1/intergroup_range.pdf}
\caption{The events and the pieces by the gap. The searches ending in a climb, the glide speed and the height lost per km of glide at the start of each event, against the gap ahead (full) and the gap behind (dashed), in the bins \IgBins~km, the last bin holding the events with none within \IgRMax~km; solo time (blue) and group time (orange), 90\% bands, bins with at least \IgMinEvents{} events. Right: $\Hk$ of the pieces against the median gap ahead and behind of the piece, in the bins of \IgHBins~km and the bin with none. The two sides of a bin hold different events: each event is binned once by its gap ahead and once by its gap behind.}
\label{fig:ig-range}
\end{figure}

\paragraph{The range of the interaction.} We read the searches and the glides at their start, with no minimum length: \IgEventsSearchSolo{} searches and \IgEventsSpeedSolo{} glides in solo time, \IgEventsSearchGroup{} and \IgEventsSpeedGroup{} in group time. The statistic of a bin is the ratio of sums over its events, and the bands resample the pilots, the blocks and their cells as before. In solo time, a search that starts with an outsider ahead within \IgBinLow~km ends in a climb \IgSearchSoloAA\% of the time, against \IgSearchSoloBA\% when the nearest outsider is behind at that distance; the difference is \IgSearchSoloDiffA{} \IgSearchSoloDiffABand{} points, and \IgSearchSoloDiffB{} \IgSearchSoloDiffBBand{} from \IgBinLow{} to \IgBinLowTwo~km. From \IgBinLowTwo{} to \IgEdgeFive~km it reverses: \IgSearchSoloDiffC{} \IgSearchSoloDiffCBand{} at \IgBinLowTwo{} to \IgEdgeThree~km, and the band of the difference includes 0 in the bin from \IgMeetSearchSolo~km on. In group time the pattern is the same, \IgSearchGroupDiffA{} \IgSearchGroupDiffABand{} points under \IgBinLow~km and \IgSearchGroupDiffD{} \IgSearchGroupDiffDBand{} at \IgEdgeThree{} to \IgEdgeFour~km; the band of the difference excludes 0 again from \IgEdgeFive{} to \IgEdgeSix~km (\IgSearchGroupDiffF{} \IgSearchGroupDiffFBand). In the bin from \IgBinLow{} to \IgBinLowTwo~km the glide speed differs by \IgSpeedSoloDiffB~km/h and the height lost per km by \IgSinkSoloDiffB~m/km between the sides in solo time, the largest differences of those two curves. For $\Hk$ by the median gap of the piece, the difference between ahead and behind in solo time is \IgHlocSoloDiffA{} \IgHlocSoloDiffABand{} in the closest bin and \IgHlocSoloDiffC{} \IgHlocSoloDiffCBand{} at \IgHEdgeTwo{} to \IgHEdgeThree~km: small, and of the sign of $\Delta_\text{info,solo}$.
So the one effect that gliders ahead have and gliders behind do not is a few points of success of the searches within about 1~km, and it vanishes, then reverses, beyond; for the straightness we see no range at which the ahead curve rises above the behind curve.

\paragraph{Limits.} The index has one position per minute: distances below a few hundred metres and delays below a minute are not resolved, so a glider in the same thermal is not told from one a few hundred metres away. Only gliders that declared their flight to the federation are in the archive: ``alone'' means no recorded glider around, and the gaps are upper bounds of the true ones. A piece changes neighbourhood along its length: the classes are shares of its time, and the lags reach 1000~s. The direction is a choice of ours (a window of \IgTDir~min, a minimum of \IgMinDisp~m), and a half-plane is the simplest definition of ahead; a cone would leave out the gliders at the side. A group has no direction when it hardly moves, and those minutes leave the shares. Being near others is itself selected, because it happens on good days, at busy sites and on frequented routes: the control behind and the within-flight line are the two guards, and neither is a randomised comparison. The within-flight line rests on far fewer flights than the raw one (Table~\ref{tab:intergroup}), and BEHIND is thin. The groups are those of one grouping run, and the time-shift chance level of the shares is missing.

\paragraph{What we keep.} The gaps allow an inter-group effect: the leader of a group has an outsider ahead within \IgRadOne, \IgRMain{} and \IgRadHigh~km during \IgShareGroupOne, \IgShareGroupFive{} and \IgShareGroupTen\% of the group time, at a median gap of \IgGapMed~km when there is one, and a member of another group within \IgRMain~km during \IgShareOtherFive\%. We have no sight range, so we cannot say how often the gliders ahead are seen. Against the intra-group effect, an inter-group effect on the straightness is not there: $\Delta_\text{intra}$ is \IgConFiveIntraRaw{} \IgConFiveIntraRawBand{} raw (\IgConFiveIntraFlight{} \IgConFiveIntraFlightBand{} within the flight), while $\Delta_\text{inter}$ is \IgConFiveInterRaw{} raw and \IgConFiveInterFlight{} \IgConFiveInterFlightBand{} within the flight. The intra-group effect is small and its within-flight band includes 0, so what we can say is that it is the only one of the two with a positive raw sign. On the phases the inter-group effect does not vanish: within the flight, gliders ahead go with more time in climb (\IgPhFiveInterClimbFlight{} points) and the group with a similar amount (\IgPhFiveIntraClimbFlight{} points).
The ahead-behind contrast does not point to information from the gliders ahead. $\Delta_\text{info,solo}$ is \IgConFiveInfoSoloRaw{} raw and \IgConFiveInfoSoloFlight{} within the flight, of the opposite sign to information, and a frequented route would give a positive $\Delta_\text{crowd}$ equal to $\Delta_\text{inter}$, which we do not see for $\Hk$ (\IgConFiveCrowdRaw{} raw, \IgConFiveCrowdFlight{} within the flight, a band that includes 0). The searches give the only hint of information, \IgSearchSoloDiffA{} points within \IgBinLow~km, with the opposite sign from \IgBinLowTwo{} to \IgEdgeFive~km. The data do not decide between information and a frequented route for the phases, where the control behind gives about as much as the gliders ahead; for the straightness they decide against both.
